\documentclass[a4paper]{article}
\usepackage{graphicx}
\usepackage{amsmath,amssymb}
\usepackage{hyperref}
\usepackage{minted}
\usepackage{multirow}
\usepackage{float}

\title{Pronominaladverb}
\date{}

\begin{document}
    \maketitle
    
    
    \section{Definition}
        
        Pronominaladverbien sind Wörter wie \textit{daran, darauf, darüber, damit, dafür}.
        Sie bestehen aus zwei Teilen:
        
        \begin{itemize}
            \item \textbf{da(r)-}
            \item einer \textbf{Präposition}
        \end{itemize}
        
        Beispiele:
        
        \begin{itemize}
            \item daran = da + an
            \item darauf = da + auf
            \item darüber = da + über
            \item damit = da + mit
            \item dafür = da + für
        \end{itemize}
        
        Wenn die Präposition mit einem Vokal beginnt, wird ein \textbf{r} eingefügt:
        
        \begin{itemize}
            \item da + an → daran
            \item da + auf → darauf
            \item da + über → darüber
        \end{itemize}
    
    
    \section{Verwendung}
        
        Pronominaladverbien verwendet man, wenn man sich auf eine \textbf{Sache} oder einen \textbf{Sachverhalt} bezieht, nicht auf eine Person.
        
        Falsch:
        
        \begin{itemize}
            \item Ich warte auf das.
            \item Ich denke an das.
        \end{itemize}
        
        Richtig:
        
        \begin{itemize}
            \item Ich warte darauf.
            \item Ich denke daran.
        \end{itemize}
    
    
    \section{Unterschied zu Personen}
        
        Für Personen benutzt man kein Pronominaladverb.
        
        \textbf{Person:}
        
        \begin{itemize}
            \item Ich denke an ihn.
            \item Ich warte auf dich.
        \end{itemize}
        
        \textbf{Sache:}
        
        \begin{itemize}
            \item Ich denke daran.
            \item Ich warte darauf.
        \end{itemize}
    
    
    \section{Verbindung mit Nebensätzen}
        
        Pronominaladverbien werden häufig mit Nebensätzen verwendet.
        
        \begin{itemize}
            \item Ich denke daran, dass er kommt.
            \item Ich arbeite daran, eine Lösung zu finden.
            \item Es hängt davon ab, wie das System funktioniert.
        \end{itemize}
        
        Hier verbindet das Pronominaladverb das Verb mit dem folgenden Nebensatz oder Infinitivsatz.
    
    
    \section{Verwendung im wissenschaftlichen Stil}
        
        Im akademischen Schreiben sind Pronominaladverbien sehr häufig.
        
        Beispiele:
        
        \begin{itemize}
            \item Wir arbeiten daran, ein Modell zu entwickeln.
            \item Wir denken darüber nach, wie Signale Bedeutung erzeugen.
            \item Es hängt davon ab, wie das System trainiert wurde.
            \item Wir konzentrieren uns darauf, die Relation zwischen Gehirn und Signal zu analysieren.
        \end{itemize}

\end{document}